% ============================================================================
% COURS 08 — EXERCICES DE CONVERSION ÉLECTROMÉCANIQUE
% ============================================================================

\section{Exercices complémentaires}\label{sec:em-exercices}

\begin{TLExercise}{Exercice 1 — Dimensionnement d'un entraînement}
Un moteur entraîne une charge par l'intermédiaire d'un réducteur de rapport
$k=\Omega_s/\Omega_m=0{,}20$ et de rendement $\eta_r=0{,}90$. La charge exige
un couple $C_s=45\,\mathrm{N\,m}$ à la vitesse $\Omega_s=12\,\mathrm{rad\,s^{-1}}$.
\begin{enumerate}
  \item Calculer la puissance mécanique utile demandée par la charge.
  \item Déterminer la vitesse du moteur.
  \item Déterminer le couple que doit fournir le moteur en régime établi.
  \item Un moteur de puissance utile nominale $700\,\mathrm W$ convient-il ? Justifier.
\end{enumerate}
\end{TLExercise}

\begin{TLExercise}{Exercice 2 — Accélération d'un axe}
Un axe ramené sur l'arbre moteur possède un moment d'inertie équivalent
$J_{eq}=0{,}18\,\mathrm{kg\,m^2}$. Le couple résistant, supposé constant, vaut
$C_r=5\,\mathrm{N\,m}$. On souhaite passer de $0$ à
$120\,\mathrm{rad\,s^{-1}}$ en $1{,}5\,\mathrm s$ avec une accélération constante.
\begin{enumerate}
  \item Calculer l'accélération angulaire.
  \item En déduire le couple utile minimal du moteur pendant l'accélération.
  \item Calculer la puissance mécanique instantanée en fin d'accélération.
  \item Comparer le couple nécessaire en accélération et en régime établi.
\end{enumerate}
\end{TLExercise}

\begin{TLExercise}{Exercice 3 — Point de fonctionnement d'un ventilateur}
La caractéristique utile d'un moteur est approchée par
$C_u=18-0{,}030\Omega$ et le ventilateur impose
$C_r=8{,}0\times10^{-4}\Omega^2$, avec les couples en newton-mètre et
$\Omega$ en radian par seconde.
\begin{enumerate}
  \item Écrire l'équation du point de fonctionnement.
  \item Déterminer numériquement la vitesse de fonctionnement positive.
  \item Calculer le couple et la puissance mécanique correspondants.
  \item Expliquer qualitativement pourquoi ce point est stable.
\end{enumerate}
\end{TLExercise}

\begin{TLExercise}{Exercice 4 — Freinage et quadrants}
Un véhicule électrique se déplace en marche avant, donc $\Omega>0$.
Lors d'un freinage récupératif, la machine exerce un couple opposé au mouvement.
\begin{enumerate}
  \item Donner le signe du couple de la machine.
  \item Donner le signe de $P_m=C\Omega$.
  \item Identifier le quadrant de fonctionnement.
  \item Indiquer les conditions nécessaires, dans toute la chaîne d'énergie, pour que
  l'énergie puisse être renvoyée vers la batterie.
  \item Que devient l'énergie si la source n'est pas réversible ?
\end{enumerate}
\end{TLExercise}

\begin{TLExercise}{Exercice 5 — Essais et rendement}
Lors d'un essai en charge, une machine absorbe $P_a=2{,}40\,\mathrm{kW}$ et fournit
un couple utile $C_u=12{,}5\,\mathrm{N\,m}$ à
$\Omega=160\,\mathrm{rad\,s^{-1}}$.
\begin{enumerate}
  \item Calculer la puissance utile.
  \item Calculer le rendement.
  \item Déterminer les pertes totales.
  \item Expliquer ce que permettrait d'estimer un essai à vide réalisé à la même vitesse.
\end{enumerate}
\end{TLExercise}

\begin{TLExercise}{Exercice 6 — Cycle de levage}
Un treuil élève puis redescend une charge. On suppose que le sens positif de vitesse
correspond à la montée et que le couple de pesanteur ramené à l'arbre est constant.
\begin{enumerate}
  \item Pour la montée à vitesse constante, donner les signes de $\Omega$, $C_m$ et $P_m$.
  \item Pour la descente freinée à vitesse constante, donner ces mêmes signes.
  \item Identifier les quadrants utilisés.
  \item Déduire si une chaîne un quadrant suffit pour cette application.
  \item Proposer une solution pour gérer l'énergie pendant la descente.
\end{enumerate}
\end{TLExercise}
